\documentclass{article}
\usepackage{amsmath,amssymb,amsthm,hyperref}
\newtheorem{theorem}{Theorem}
\begin{document}
\title{Low-rank tensor residues and compact exterior oscillation}
\maketitle

\paragraph{Status.}
The scalar identities and continuity argument below are complete written
proofs.  Their displayed endpoint inequalities have an independent exact
rational certificate.  They are not Lean declarations.  No growing local
decoder, admissible tensor correction, or new asymptotic bound for the
original worst-case quantity follows from this note.

\section{The exact tensor family}

Fix $c$ pole blocks $p_l(0),\ldots,p_l(q-1)$.  A residue tensor of CP rank
at most $R$ has a decomposition
\[
 A(\alpha_1,\ldots,\alpha_c)
 =\sum_{r=1}^R\prod_{l=1}^c a_{rl}(\alpha_l).
\]
Its scalar exterior Cauchy kernel is
\[
 F(z)=\sum_{\alpha\in[q]^c}
 \frac{A(\alpha)}{z+\sum_l p_l(\alpha_l)},
 \qquad z> -\min_\alpha\sum_l p_l(\alpha_l).
\]
After shifting each pole block, suppose all poles are nonnegative.  Then
\[
 F(z)=\int_0^\infty e^{-zu}G(u)\,du,
 \qquad G(u)=\sum_{r=1}^R\prod_{l=1}^c E_{rl}(u),
 \qquad E_{rl}(u)=\sum_a a_{rl}(a)e^{-p_l(a)u}.
\]
For $R=1$, each nonzero exponential polynomial $E_{1l}$ has at most
$q-1$ real zeros.  Thus $G$ has at most $c(q-1)$ sign changes.  The
variation-diminishing property of the Laplace kernel gives the same
upper bound for the number of distinct zeros of nonzero $F$ on its
exterior half-line.  This argument does not extend by simply adding
the root counts of two product terms.

For integer poles, $x=e^{-u}$ identifies $G$ with a sum of $R$ products
of sparse polynomials.  The common pole support across the $R$ terms
is not a substantial restriction: for two products of $t$-sparse
polynomials, pad each pair of factor supports to their union, of size
at most $2t$, and use zero coefficients on missing exponents.  After
rescaling the polynomial variable, any finite collection of positive
roots lies in $(0,1)$.  Therefore a polynomial bound in $(c,q)$ for
zeros of all rank-two densities $G$ would imply the two-summand case
of the real tau-conjecture for positive roots.

Grenet, Koiran, Portier and Strozecki explicitly identify the
two-summand case as unresolved in their paper, and the Wronskian
bound of Koiran, Portier and Tavenas is exponential in the number of
distinct sparse factors.  This is a precise connection for the
\emph{density} family.  It does not assert equivalence between the
real tau-conjecture and the number of exterior zeros of its Laplace
transform.  Variation diminution gives only the upper-bound direction.

\section{Rank two can attain all eight available exterior zeros}

\begin{theorem}
There are two three-element nonnegative pole blocks and a real residue
matrix of rank exactly two whose scalar exterior kernel has eight
distinct positive zeros in an arbitrarily short interval.  Every zero
lies strictly beyond the full-cube gap budget for the generic
three-point-plane extension described in the tensor audit.
\end{theorem}
\begin{proof}
Use the fixed blocks
\[
 p=(0,2947,9807),\qquad q=(0,202,1015).
\]
All nine sums $\lambda_{ab}=p_a+q_b$ are distinct.  Their optimized
generic-row extension has the gap budget
\[
 B=(3\cdot9807-2\cdot2947)+(3\cdot1015-2\cdot202)=26168.
\]
For $D>B$ and $\varepsilon\geq0$, put
\[
 P_{D,\varepsilon}(z)=\prod_{i=0}^7(z-D-i\varepsilon),
 \qquad Q(z)=\prod_{a,b}(z+\lambda_{ab}),
 \qquad
 A_{ab}(D,\varepsilon)=
 \frac{P_{D,\varepsilon}(-\lambda_{ab})}
 {\prod_{(a',b')\ne(a,b)}(\lambda_{ab}-\lambda_{a'b'})}.
\]
The denominator in the last expression equals $Q'(-\lambda_{ab})$
because there are eight factors.  Equality of residues and decay at
infinity prove the rational identity
\[
 \sum_{a,b}\frac{A_{ab}(D,\varepsilon)}{z+\lambda_{ab}}
 =\frac{P_{D,\varepsilon}(z)}{Q(z)}.
 \tag{1}
\]
Let $D_-=1001B/1000=3274271/125$ and $D_+=2B=52336$.
The exact rational checker verifies
\[
 \det A(D_-,0)<0<\det A(D_+,0).
 \tag{2}
\]
It also verifies these same two strict signs at $\varepsilon=1/10$.
The complete rational values, rather than decimal approximations,
are saved in the JSON certificate.

Each determinant is a polynomial in $(D,\varepsilon)$ times a fixed
nonzero rational constant.  Continuity of (2) therefore supplies
$\varepsilon_0>0$ such that its strict signs persist whenever
$0\leq\varepsilon<\varepsilon_0$.  For each positive such
$\varepsilon$, the intermediate value theorem supplies
$D_*\in(D_-,D_+)$ with $\det A(D_*,\varepsilon)=0$.
Thus the matrix has rank at most two.  Identity (1) gives the eight
distinct zeros $D_*+i\varepsilon$, all in an interval of width
$7\varepsilon$.  They are all beyond $B$, with the uniform gap
$D_--B=3271/125$.  The width can be made arbitrarily small.

The matrix is nonzero by (1).  It cannot have rank one, because the
rank-one Laplace product argument gives at most $2(3-1)=4$ exterior
zeros.  Its rank is consequently exactly two.
\end{proof}

The gap and arbitrarily small interval leave room for any fixed finite
column-token budget and strict negative filler slopes.  This statement
checks the scalar denominator budget only.  Implementing the rank-two
residues and repairing all tensor partial diagonals still requires the
separate admissible-score construction; that construction is not claimed
here.  Eight zeros at $(c,q,R)=(2,3,2)$ do not imply exponentially many
zeros at growing $c$.

The independent verifier is
\begin{verbatim}
python3 artifacts/head-compression/JointCellRankTwoExteriorCertificate.py
\end{verbatim}
It uses only integers and rational arithmetic.  It verifies both endpoint
sign pairs, distinct poles, the gap budget, and rational partial-fraction
evaluations.  The proof of (1) and the continuity argument establish the
uniform statements beyond those finite evaluations.

\section{What tensor trains would additionally require}

A bounded-width matrix-product residue can have exponentially many CP
terms, so it is not covered by a small-$R$ product-density argument.
However, one mixed derivative of a scalar affine reciprocal gives a
product of one factor from each row block.  A sum of $R$ such scalar
channels gives precisely $R$ CP terms.  A noncommuting matrix-product
description by itself does not supply an admissible scalar head.

Encoding adjacent block pairs as new one-hot coordinates makes their
local factors affine, but forces consistency between neighboring pairs.
For $q^2$ states per original block, the consistent tuples number
$q^{2c}$ inside a naive $q^{4c}$ independent pair-coordinate domain.
Using only those tuples loses an exponential factor in that direct
cover.  A useful tensor-train route therefore needs a new exact score or
a cover that charges these consistency constraints without this loss.
No impossibility statement for more general realizations is asserted.

\paragraph{Primary sources.}
\href{https://www-verimag.imag.fr/~grenetb/publis/GKPS11.pdf}
{Grenet, Koiran, Portier and Strozecki, The Limited Power of Powering};
\href{https://arxiv.org/pdf/1205.1015}
{Koiran, Portier and Tavenas, A Wronskian Approach to the Real
\(\tau\)-Conjecture}.
\end{document}
